% Restriction, Range and Zigzag


%restriction
\newcommand{\rst}[1]{\overline {\vphantom{f}#1}}
\newcommand{\rstItself}{$\overline{\kern0.2cm\nudgeup{0.2cm}}$}

% range  
\usepackage{scalerel,stackengine}
\stackMath
\newcommand\reallywidehat[1]{%
\savestack{\tmpbox}{\stretchto{%
  \scaleto{%
    \scalerel*[\widthof{\ensuremath{#1}}]{\kern-.6pt\bigwedge\kern-.6pt}%
    {\rule[-\textheight/2]{1ex}{\textheight}}%WIDTH-LIMITED BIG WEDGE
  }{\textheight}% 
}{0.5ex}}%
\stackon[1pt]{#1}{\tmpbox}%
}
\parskip 1ex

\newcommand{\rng}[1]{\reallywidehat {\vphantom{f}#1}}
\newcommand{\rngItself}{$\widehat{\kern0.2cm\nudgeup{0.2cm}}$}

%Experimental
\newcommand{\weakrange}[1]{\widehat{#1}^{\mathrm w}}
\NewDocumentCommand{\pointwiserange}{m}
  {\reallywidehat{#1}^{\!*}}

%Stickleback category -- will probably rename this
\newcommand{\Stk}[1][C]{\textbf{Stk}_\textbf{#1}}

%pHom that didn't work out 
% Only need this if I am talking about what doesn't work.
\NewDocumentCommand{\HomStk}{O{a} O{-} O{C}}{pHom_{\Stk[#3]}(#1,#2)}

%%%%%%%%%%%%%%%%%%%%%%%%%%%%%%%%%%%%%
% Ob macros used in pstricks diagrams
%    might these be used more widely?
%%%%%%%%%%%%%%%%%%%%%%%%%%%%%%%%%%%%%

\newcommand{\Ob}[2]{
\Rnode{#1#2}{#1_{#2}}
}
% b bare object no subsript
\newcommand{\bOb}[1]{
\Rnode{#1}{#1}
}
% p for predecessor
\newcommand{\Obp}[2]{
\Rnode{#1#2p}{#1_{#2-1}}
}

% s for successor
\newcommand{\Obs}[2]{
\Rnode{#1#2s}{#1_{#2+1}}
}

\newcommand{\Obpp}[2]{
\Rnode{#1#2pp}{#1_{#2-2}}
}

\newcommand{\Obss}[2]{
\Rnode{#1#2ss}{#1_{#2+2}}
}
%%%%%%%%%%%%%%%%%%%%%%%
% zigzags as tuples
%%%%%%%%%%%%%%%%%%%%%%%

%zigzagTuple{f}{n}
\newcommand{\zigzagTuple}[2]
{
\tuple{#1_1,#1'_1,...#1_{#2-1},#1'_{#2-1},#1_{#2}}    
}

%zigzagTupleLedIn{f}{g}{m}    GIVES <f \circ g_1,...g_m>
\newcommand{\zigzagTupleLedIn}[3]
{
\tuple{#1 \circ #2_1,#2'_1,...#1_{#3-1},#2'_{#3-1},#2_{#3}}   
}

%zigzagTupleLedOut{f}{n}{g}    GIVES <f_1,...f_n \circ g>
\newcommand{\zigzagTupleLedOut}[3]
{
\tuple{#1_1,#1'_1,...#1_{#2-1},#1'_{#2-1},#1_{#2} \circ #3}    
}

%zigzagLeadingTuple{f}{n} GIVES <f_1,...f_n-1>
\newcommand{\zigzagLeadingTuple}[2]
{
\tuple{#1_1,#1'_1,...#1_{#2-2},#1'_{#2-2},#1_{#2-1}}    
}

%zigzagStrippedHeadTuple{f}{n}
\newcommand{\zigzagBareHeadTuple}[2]
{
\tuple{\rng{#1_1},#1'_1,...#1_{#2-1},#1'_{#2-1},#1_{#2}}    
}

%zigzagStrippedTuple{f}{n}
\newcommand{\zigzagBareTailTuple}[2]
{
\tuple{#1_1,#1'_1,...#1_{#2-1},#1'_{#2-1},\rst{#1_{#2}}}     
}

%zigzagTrailingTuple{f}{n}
\newcommand{\zigzagTrailingTuple}[2]
{
\tuple{#1_2,#1'_2,...#1_{#2-1},#1'_{#2-1},#1_{#2}}     
}

%zigzagTupleSharp{f}{n}
\newcommand{\zigzagTupleSharp}[2]
{
\tuple{#1'_1,...#1_{#2-1},#1'_{#2-1},#1_{#2}}    
}

%zigzagTupleEndShortened{f}{n}{s}
\newcommand{\zigzagTupleEndShortened}[3]
{
\tuple{#1_1,#1'_1,...#1'_{#2-1}, #1_{#2-1} \circ s}    
}

%zigzagTailTuple{f}{n}{m}
\newcommand{\zigzagTailTuple}[3]
{
\tuple{#1_#2,#1'_#2,...#1_{#2+#3-1},#1'_{#2+#3-1},#1_{#2+#3}}    
}

\newcommand{\restrictToDSource}[2]{
\overline{D(#1_1)}\circ ... \circ \overline{D(#1_#2)}
}

\newcommand{\restrictToSource}[2]{
\overline{#1_1}\circ ... \circ \overline{#1_#2}
}

\newcommand{\conc}{::}

%%%%%%%%%%%%%%%%%%%%%%%%%%%%%%%%%%%%%%%%%%%%%%%
% zigzag diagrams and zigzag diagram fragments
%%%%%%%%%%%%%%%%%%%%%%%%%%%%%%%%%%%%%%%%%%%%%%%

%\zigzagTwo{x}{y}{f}[g] %g is optional for circ f1
% This is a zig zag zig zag. 
% zig one is  f1 : x1 -> y1 
% zig one is  f'1 : x2 -> y1
\NewDocumentCommand{\zigzagTwo}{m m m o}
{
\begin{array} {c c c c c}
 &\Ob{#2}{1}&&\Ob{#2}{2}          &\\[0.8cm]
\IfNoValueTF{#4}{\Ob{#1}{1}}{\Rnode{#11}{}}&&\Ob{#1}{2}&&\Ob{#1}{3} 
\end{array}
\begin{arrows}
\ncarr{#11}{#21} \alabel{ 
\IfNoValueTF{#4}
    {}
    {#4 \circ} #3_1
}[0.45][0.4]
\ncarr{#12}{#21} \blabel{#3'_1}[0.45][0.4]
\ncarr{#12}{#22} \alabel{#3_2}[0.45][0.4]
\ncarr{#13}{#22} \blabel{#3'_2}[0.45][0.4]
\end{arrows}
}

\NewDocumentCommand{\zigzagOneAndAHalf}{m m m o}{
\begin{array} {c c c c }
 &\Ob{#2}{1}&&\Ob{#2}{2}          \\[0.8cm]
\Ob{#1}{1}&&\Ob{#1}{2}&
\end{array}
\begin{arrows}
\ncarr{#11}{#21} \alabel{#3_1 
\IfNoValueTF{#4}
    {}
    {\circ #4}
}[0.45][0.4]
\ncarr{#12}{#22} \alabel{#3_2}[0.45][0.4]
\ncarr{#12}{#21} \blabel{#3'_1}[0.45][0.4]
\end{arrows}
}

\NewDocumentCommand{\zigzagOne}{m m m o}{
\begin{array} {c c c c }
 &\Ob{#2}{1}&&          \\[0.8cm]
\Ob{#1}{1}&&\Ob{#1}{2}&
\end{array}
\begin{arrows}
\ncarr{#11}{#21} \alabel{#3_1 
\IfNoValueTF{#4}
    {}
    {\circ #4}
}[0.45][0.4]
\ncarr{#12}{#21} \blabel{#3'_1}[0.45][0.4]
\end{arrows}
}

%\zigzagZigiZag{x}{y}{f}[i]
% starts from i with a zig and a zag 
\NewDocumentCommand{\zigzagZigiZag}{m m m m}{
\begin{array} {c c c c }
 &\Ob{#2}{#4}&&\Obs{#2}{#4}    \\[0.8cm]
\Ob{#1}{#4}&&\Obs{#1}{#4}& 
\end{array}
\begin{arrows}
\ncarr{#1#4}{#2#4} \alabel{#3_{#4} }[0.6][0.4]
\ncarr{#1#4s}{#2#4} \blabel{#3'_{#4} }[0.6][0.4]
\ncarr{#1#4s}{#2#4s} \blabel{#3_{#4+1} }[0.7][0.4]
\end{arrows}
}

%\zigzagZigiZagShortening{x}{y}{f}[i][s]
\NewDocumentCommand{\zigzagZigiZagShortening}{m m m m m}
{
\zigzagZigiZag{#1}{#2}{#3}{#4}
\begin{arrows}
\ncarr[15]{#2#4}{#2#4s} \alabel{#5}
\end{arrows}
}

%\zigzagZigiZagShortened{x}{y}{f}[i]
% starts from i with a zig and a zag 
\NewDocumentCommand{\zigzagZigiZagShortened}{m m m m m}
{
\begin{array} {c c c c}
\Obp{#2}{#4}&            &\Obs{#2}{#4}&   \\[0.8cm]
            & \Ob{#1}{#4}&            &\Obss{#1}{#4} 
\end{array}
\begin{arrows}
\ncarr{#1#4}{#2#4p} \alabel{#3'_{#4-1} }[0.6][0.2]
\ncarr{#1#4}{#2#4s} \alabel{#3_{#4} \circ #5 }[0.6][0.4]
\ncarr{#1#4ss}{#2#4s} \blabel{#3'_{#4+1} }[0.5][0.3]
\end{arrows}
}


%\zigzagStart{x}{y}{f}[n]
% arrows starts from n instead of one 
\NewDocumentCommand{\zigzagStart}{m m m m}{
\begin{array} {c c c c  }
 &\Ob{#2}{1}&&\Ob{#2}{2}    \\[0.8cm]
\Ob{#1}{1}&&\Ob{#1}{2}& 
\end{array}
\begin{arrows}
\ncarr{#11}{#21} \alabel{#3_{#4} }[0.6][0.4]
\ncarr{#12}{#21} \blabel{#3'_{#4} }[0.6][0.4]
\ncarr{#12}{#22} \blabel{#3_{#4+1} }[0.7][0.4]
\end{arrows}
}


%\zigzagDotsThenFinal{x}{y}{f}{n}
% This is a zig
% fn  : yn -> xn
\NewDocumentCommand{\zigzagDotsThenFinal}{m m m m}{
\begin{array}{c c c}
\dots &&\Ob{#2}{#4} \\[0.8cm]
\dots &\Ob{#1}{#4}
\end{array}
\begin{arrows}
\ncarr{#1#4}{#2#4} \alabel{#3_{#4}}[0.6][0.4]
\end{arrows}    
}

%\zigzagDotsThenLongFinal{x}{y}{f}{n}
% This is a zig
% fn  : yn -> xn
\NewDocumentCommand{\zigzagDotsThenLongFinal}{m m m m}{
\begin{array}{c c c c c}
\dots &&\Obp{#2}{#4}&&\Ob{#2}{#4} \\[0.8cm]
\dots &\Obp{#1}{#4}&&\Ob{#1}{#4}
\end{array}
\begin{arrows}
\ncarr{#1#4p}{#2#4p} \alabel{#3_{#4-1}}[0.5][0.4]
\ncarr{#1#4}{#2#4p} \blabel{#3'_{#4-1}}[0.5][0.4]
\ncarr{#1#4}{#2#4} \blabel{#3_{#4}}[0.6][0.4]
\end{arrows}    
}

%\zigzagDotsThenFinalFlex{x}{y}{f}{n}{m}
% This is zig
% fm : yn -> xn
\NewDocumentCommand{\zigzagDotsThenFinalFlex}{m m m m m}{
\begin{array}{c c c}
\dots &&\Ob{#2}{#4} \\[0.8cm]
\dots &\Ob{#1}{#4}
\end{array}
\begin{arrows}
\ncarr{#1#4}{#2#4} \alabel{#3_{#5}}[0.6][0.4]
\end{arrows}    
}

%\zigzag{x}{y}{f}[ff]{n}
% This is a zig zag zig zag... zig
% For params 1 to 4, see command \zigzagTwo
% For param 5 see command \zigzagDotsThenFinal
%\zigzag{x}{y}{f}{n}
\NewDocumentCommand{\zigzag}{m m m o m}{
\zigzagTwo{#1}{#2}{#3}[#4]
\zigzagDotsThenFinal{#1}{#2}{#3}{#5}
}

%\zigzagFlex{x}{y}{f}{n}{n+m}
\NewDocumentCommand{\zigzagFlex}{m m m m m}{
\zigzagStart{#1}{#2}{#3}{#4}
\zigzagDotsThenFinalFlex{#1}{#2}{#3}{#4}{#5}
}

\newcommand{\zigzagfgcomposite} %Maybe this becomes zigzagComposite?
{
\arraycolsep=7pt
\zigzagTwo{w}{x}{f}
\begin{array}{c}
\dots \\[0.8cm]
\dots
\end{array}
\begin{array}{c c}
\Obp{x}{n} \\[0.8cm]
& \Ob{w}{n}
\end{array}
\begin{arrows}
\ncarr{wn}{xnp} \alabel{f'_{n-1}}[0.7][0.4] %0.2 was 0.45
\end{arrows}
\kern-15pt
\zigzag{y}{z}{g}[f_n]{m}
}

%\zigzagShortenable{x}{y}{f}{n}{i}{s}
\NewDocumentCommand{\zigzagShortenable}{m m m m m m}
{
\zigzagOne{#1}{#2}{#3}
\begin{array}{c}
\dots \\[0.8cm]
\dots
\end{array}
\zigzagZigiZagShortening{#1}{#2}{#3}{#5}{#6}
\zigzagDotsThenFinal{x}{y}{f}{n}
}

%\zigzagShortened{x}{y}{f}{n}{i}{s}
\NewDocumentCommand{\zigzagShortened}{m m m m m m}
{
\zigzagOne{#1}{#2}{#3}
\begin{array}{c}
\dots \\[0.8cm]
\dots
\end{array}
\zigzagZigiZagShortened{#1}{#2}{#3}{#5}{#6}
\zigzagDotsThenFinal{x}{y}{f}{n}
}

%\zigzagSequentialShortening{x}{y}{f}{n}{x}{g}{h}{s}
\NewDocumentCommand{\zigzagSequentialShortening}
                    {m m m m m m m m }
{
\begin{array} {c}
\\[0.75cm]
\zigzagOneAndAHalf{#1}{#2}{#3}
\zigzagDotsThenLongFinal{#1}{#2}{#3}{#4}
\end{array}
\begin{arrows}
\ncarrarc{#21}{#2#4}{60}{80} \alabel{#8_1}[0.3]
\ncarrarc{#22}{#2#4}{45}{40} \alabel{#8_2}[0.3]
\ncarrarc{#2#4p}{#2#4}{10}{10} \alabel{#8_{#4-1}}[0.2]
\end{arrows}
}
